\begin{register}{H}{IO\_MUX\_PIN\_CTRL}{0x3FF49000}\label{regdesc:IOMUXPINCTRL}
\regfield{(reserved)}{20}{12}{{0x0}}%
\regfield{PIN\_CTRL\_CLK3}{4}{8}{{0x0}}%
\regfield{PIN\_CTRL\_CLK2}{4}{4}{{0x0}}%
\regfield{PIN\_CTRL\_CLK1}{4}{0}{{0x0}}%
\reglabel{Reset}\regnewline%
\begin{regdesc}
{\bfseries 要将 I2S\regindex{0} 外设时钟 (I2S\regindex{0}\_CLK) 输出到：}\\
CLK\_OUT1，配置 PIN\_CTRL[3:0] = 0x0；\\
CLK\_OUT2，配置 PIN\_CTRL[3:0] = 0x0 且 PIN\_CTRL[7:4] = 0x0；\\
CLK\_OUT3，配置 PIN\_CTRL[3:0] = 0x0 且 PIN\_CTRL[11:8] = 0x0。\\
{\bfseries 要将 I2S\regindex{1} 外设时钟 (I2S\regindex{1}\_CLK) 输出到：}\\
CLK\_OUT1，配置 PIN\_CTRL[3:0] = 0xF；\\
CLK\_OUT2，配置 PIN\_CTRL[3:0] = 0xF 且 PIN\_CTRL[7:4] = 0x0；\\
CLK\_OUT3，配置 PIN\_CTRL[3:0] = 0xF 且 PIN\_CTRL[11:8] = 0x0。

{\bfseries 要将 APLL 输出到}\\
CLK\_OUT1，配置 PIN\_CTRL[3:0] = 0x6；\\
CLK\_OUT2，配置 PIN\_CTRL[3:0] = 0x6 且 PIN\_CTRL[7:4] = 0x6；\\
CLK\_OUT3，配置 PIN\_CTRL[3:0] = 0x6 且 PIN\_CTRL[11:8] = 0x6。（读/写）
\end{regdesc}
\end{register}

\vspace{-2em}
\begin{tiplisting}
\vspace{-2em}
\begin{itemize}
\item 仅支持上述配置组合，即；
\begin{itemize}
    \item 仅支持输出时钟：I2S\regindex{0/1}\_CLK 和 APLL;
    \item 仅支持输出管脚：CLK\_OUT1-3
\end{itemize}

\item CLK\_OUT1-3 可在 \hyperref[table:iomux_pads]{IO\_MUX Pad 列表}中查询。
\end{itemize}
\end{tiplisting}


\begin{register}{H}{IO\_MUX\_\textcolor{red}{\it{x}}\_REG (\textcolor{red}{\it{x}}: 0-19, 20$^1$, 21-23, 25-27, 32-39)}{地址见表 \ref{tab:iomux_reg_list}}

\label{regdesc:IOMUXnREG}
\regfield{(reserved)}{17}{15}{00000000000000000}%
\regfield{MCU\_SEL}{3}{12}{{0x0}}%
\regfield{FUN\_DRV}{2}{10}{{0x2}}%
\regfield{FUN\_IE}{1}{9}{0}%
\regfield{FUN\_WPU}{1}{8}{0}%
\regfield{FUN\_WPD}{1}{7}{0}%
\regfield{MCU\_DRV}{2}{5}{{0x0}}%
\regfield{MCU\_IE}{1}{4}{0}%
\regfield{MCU\_WPU}{1}{3}{0}%
\regfield{MCU\_WPD}{1}{2}{0}%
\regfield{SLP\_SEL}{1}{1}{0}%
\regfield{MCU\_OE}{1}{0}{0}%
\reglabel{Reset}\regnewline%
\begin{regdesc}\begin{reglist}
\item [MCU\_SEL]为信号选择 IO\_MUX 功能。0：选择功能 0；1：选择功能 1；以此类推。（读/写）
\label{FUNDRV} \item [FUN\_DRV]选择驱动强度，值越大，强度越大。对于 GPIO34-39，FUN\_DRV 恒为 0。具体驱动强度请参考 \linkchipdatasheetCN 中“管脚清单说明”的第 8 条。（读/写）
\item [FUN\_IE] 管脚的输入使能。1：输入使能；0：输入关闭。（读/写）
\item [FUN\_WPU] \label{regdesc:FUNWPU} 管脚的上拉使能。1：内部上拉使能；0：内部上拉关闭。由于 GPIO34-39 仅可用作输入管脚，不具备输出驱动或内置上拉/下拉电路，因此这些管脚的 FUN\_WPU 恒为 0。（读/写）
\item [FUN\_WPD] \label{regdesc:FUNWPD} 管脚的下拉使能。1：内部下拉使能；0：内部下拉关闭。由于 GPIO34-39 仅可用作输入，不具备输出驱动或内置上拉/下拉电路，因此这些管脚的 FUN\_WPD 恒为 0。（读/写）
\item [MCU\_DRV]睡眠模式下选择管脚的驱动强度，值越大，强度越大。（读/写）
\item [MCU\_IE]睡眠模式下管脚的输入使能。1：输入使能；0：输入关闭。（读/写）
\item [MCU\_WPU]睡眠模式下管脚的上拉使能。1：内部上拉使能；0：内部上拉关闭。（读/写）
\item [MCU\_WPD]睡眠模式下管脚的下拉使能。1：内部下拉使能；0：内部下拉关闭。（读/写）
\item [SLP\_SEL]管脚的睡眠模式选择。置 1 将使能睡眠模式。（读/写）
\item [MCU\_OE]睡眠模式下管脚的输出使能。1:输出使能；2:输出关闭。（读/写）
\end{reglist}\end{regdesc}

\begin{tiplisting}
\vspace{-1em}
\begin{enumerate}
  \vspace{-1em}
  \item GPIO20 仅在 ESP32-PICO-V3 和 ESP32-PICO-V3-02 中可用。
\end{enumerate}
\end{tiplisting}
\end{register}